% GENERATED FILE. Edit canonical lessons, then run npm run generate:books.
% canonical-source-hash: fnv1a64:db46ebf208209e49
% canonical-lessons: RU-C513-razminka, RU-C513-otzhimatsya, RU-C513-prisedat, RU-C513-ganteli, RU-C513-vynoslivost

\chapter{Warming Up}
\label{ch:ru-warming-up}
\hlchaptermodality{513}

\begin{chapteropening}
\textbf{By the end of this chapter:} \emph{I can talk about exercises in the gym.}

Complete the last lesson of chapter 513: I can talk about exercises in the gym.
\end{chapteropening}

\section[\texorpdfstring{\ru{разминка} (razmínka)}{разминка (razmínka)}]{\ru{разминка} (razmínka) — a warm-up}

\label{lesson:RU-C513-razminka}



\begin{quote}
Before the new one: say the Russian for to fall behind, then the Russian for a distance.
\end{quote}

\subsection*{You'll want to know: \ru{разминка}}
\textbf{\ru{разминка}} (\emph{razmínka}) — \textquotedblleft{}a warm-up\textquotedblright{}.

\begin{quote}
\ru{Перед} \ru{бегом} \ru{обязательно} \ru{нужна} \ru{разминка}.
\end{quote}

\begin{quote}
\emph{Péred bégom obyazátelno nuzhná razmínka.}
\end{quote}

\begin{quote}
A warm-up is essential before a run.
\end{quote}

Feminine.

\begin{cousinweb}
\emph{\ru{разминка}} is built on \emph{\ru{размять}} 'to loosen, to knead', from \emph{\ru{мять}} 'to crumple'; EtymDB has no entry.
\end{cousinweb}

\begin{grammarlens}[title={What you've built}]
One more word for telling, explaining and describing at B1.
\end{grammarlens}

\subsection*{Your turn}
\begin{itemize}
\raggedright
  \item \emph{Say it:} \emph{razmínka}
  \item \emph{Say it:} \emph{razmínka}, once more
  \item \emph{Say it:} a sentence of your own with \emph{razmínka}
  \item \emph{Recall:} say the Russian for to fall behind, then the Russian for a distance, then say \emph{razmínka} again
\end{itemize}

\begin{tcolorbox}[breakable,colback=teal!4,colframe=teal!35!black,title={Before you move on}]
What does \ru{разминка} mean? (\textquotedblleft{}A warm-up\textquotedblright{}.) Say it once more.
\end{tcolorbox}
\section[\texorpdfstring{\ru{отжиматься} (otzhimátsya)}{отжиматься (otzhimátsya)}]{\ru{отжиматься} (otzhimátsya) — to do press-ups}

\label{lesson:RU-C513-otzhimatsya}



\begin{quote}
Before the new one: say the Russian for a distance, then the Russian for a warm-up.
\end{quote}

\subsection*{You'll want to know: \ru{отжиматься}}
\textbf{\ru{отжиматься}} (\emph{otzhimátsya}) — \textquotedblleft{}to do press-ups\textquotedblright{}.

\begin{quote}
\ru{Каждое} \ru{утро} \ru{он} \ru{отжимается} \ru{двадцать} \ru{раз}.
\end{quote}

\begin{quote}
\emph{Kázhdoye útro on otzhimáyetsya dvádtsat raz.}
\end{quote}

\begin{quote}
Every morning he does twenty press-ups.
\end{quote}

Imperfective; perfective \ru{отжаться} (\ru{отожмусь}).

\begin{cousinweb}
\emph{\ru{отжиматься}} is built on \emph{\ru{от}-} and \emph{\ru{жать}} 'to press, to squeeze'; EtymDB has no entry.
\end{cousinweb}

\begin{grammarlens}[title={What you've built}]
One more word for telling, explaining and describing at B1.
\end{grammarlens}

\subsection*{Your turn}
\begin{itemize}
\raggedright
  \item \emph{Say it:} \emph{otzhimátsya}
  \item \emph{Say it:} \emph{otzhimátsya}, once more
  \item \emph{Say it:} a sentence of your own with \emph{otzhimátsya}
  \item \emph{Recall:} say the Russian for a distance, then the Russian for a warm-up, then say \emph{otzhimátsya} again
\end{itemize}

\begin{tcolorbox}[breakable,colback=teal!4,colframe=teal!35!black,title={Before you move on}]
What does \ru{отжиматься} mean? (\textquotedblleft{}To do press-ups\textquotedblright{}.) Say it once more.
\end{tcolorbox}
\section[\texorpdfstring{\ru{приседать} (prisedát)}{приседать (prisedát)}]{\ru{приседать} (prisedát) — to squat}

\label{lesson:RU-C513-prisedat}



\begin{quote}
Before the new one: say the Russian for a warm-up, then the Russian for to do press-ups.
\end{quote}

\subsection*{You'll want to know: \ru{приседать}}
\textbf{\ru{приседать}} (\emph{prisedát}) — \textquotedblleft{}to squat\textquotedblright{}.

\begin{quote}
\ru{Тренер} \ru{сказал} \ru{нам} \ru{приседать} \ru{медленно}.
\end{quote}

\begin{quote}
\emph{Tréner skazál nam prisedát médlenno.}
\end{quote}

\begin{quote}
The coach told us to squat slowly.
\end{quote}

Imperfective; perfective \ru{присесть} (\ru{присяду}), which also means to sit down for a moment.

\begin{cousinweb}
\emph{\ru{приседать}} is built on \emph{\ru{при}-} and the root of \emph{\ru{сесть}} 'to sit down'; EtymDB lists the verb without a source.
\end{cousinweb}

\begin{grammarlens}[title={What you've built}]
One more word for telling, explaining and describing at B1.
\end{grammarlens}

\subsection*{Your turn}
\begin{itemize}
\raggedright
  \item \emph{Say it:} \emph{prisedát}
  \item \emph{Say it:} \emph{prisedát}, once more
  \item \emph{Say it:} a sentence of your own with \emph{prisedát}
  \item \emph{Recall:} say the Russian for a warm-up, then the Russian for to do press-ups, then say \emph{prisedát} again
\end{itemize}

\begin{tcolorbox}[breakable,colback=teal!4,colframe=teal!35!black,title={Before you move on}]
What does \ru{приседать} mean? (\textquotedblleft{}To squat\textquotedblright{}.) Say it once more.
\end{tcolorbox}
\section[\texorpdfstring{\ru{гантели} (gantéli)}{гантели (gantéli)}]{\ru{гантели} (gantéli) — dumbbells}

\label{lesson:RU-C513-ganteli}



\begin{quote}
Before the new one: say the Russian for to do press-ups, then the Russian for to squat.
\end{quote}

\subsection*{You'll want to know: \ru{гантели}}
\textbf{\ru{гантели}} (\emph{gantéli}) — \textquotedblleft{}dumbbells\textquotedblright{}.

\begin{quote}
\ru{Дома} \ru{у} \ru{меня} \ru{есть} \ru{гантели} \ru{и} \ru{коврик}.
\end{quote}

\begin{quote}
\emph{Dóma u menyá yest gantéli i kóvrik.}
\end{quote}

\begin{quote}
At home I have dumbbells and a mat.
\end{quote}

Plural of \emph{\ru{гантель}} (feminine).

\begin{cousinweb}
\emph{\ru{гантели}} is the plural of \emph{\ru{гантель}}, borrowed from German \emph{Hantel} 'dumbbell' (EtymDB).
\end{cousinweb}

\begin{grammarlens}[title={What you've built}]
One more word for telling, explaining and describing at B1.
\end{grammarlens}

\subsection*{Your turn}
\begin{itemize}
\raggedright
  \item \emph{Say it:} \emph{gantéli}
  \item \emph{Say it:} \emph{gantéli}, once more
  \item \emph{Say it:} a sentence of your own with \emph{gantéli}
  \item \emph{Recall:} say the Russian for to do press-ups, then the Russian for to squat, then say \emph{gantéli} again
\end{itemize}

\begin{tcolorbox}[breakable,colback=teal!4,colframe=teal!35!black,title={Before you move on}]
What does \ru{гантели} mean? (\textquotedblleft{}Dumbbells\textquotedblright{}.) Say it once more.
\end{tcolorbox}
\section[\texorpdfstring{\ru{выносливость} (vynóslivost)}{выносливость (vynóslivost)}]{\ru{выносливость} (vynóslivost) — stamina, endurance}

\label{lesson:RU-C513-vynoslivost}



\begin{quote}
Before the new one: say the Russian for to squat, then the Russian for dumbbells.
\end{quote}

\subsection*{You'll want to know: \ru{выносливость}}
\textbf{\ru{выносливость}} (\emph{vynóslivost}) — \textquotedblleft{}stamina, endurance\textquotedblright{}.

\begin{quote}
\ru{Для} \ru{марафона} \ru{нужна} \ru{выносливость}.
\end{quote}

\begin{quote}
\emph{Dlya marafóna nuzhná vynóslivost.}
\end{quote}

\begin{quote}
You need stamina for a marathon.
\end{quote}

Feminine. Built on \emph{\ru{выносливый}} — tough, hardy.

\begin{cousinweb}
\emph{\ru{выносливость}} is built on \emph{\ru{выносливый}} 'hardy', from \emph{\ru{выносить}} 'to bear, to endure' and \emph{\ru{носить}} 'to carry'; EtymDB has no entry.
\end{cousinweb}

\begin{grammarlens}[title={What you've built}]
One more word for telling, explaining and describing at B1.
\end{grammarlens}

\subsection*{Your turn}
\begin{itemize}
\raggedright
  \item \emph{Say it:} \emph{vynóslivost}
  \item \emph{Say it:} \emph{vynóslivost}, once more
  \item \emph{Say it:} a sentence of your own with \emph{vynóslivost}
  \item \emph{Recall:} say the Russian for to squat, then the Russian for dumbbells, then say \emph{vynóslivost} again
\end{itemize}

\begin{tcolorbox}[breakable,colback=teal!4,colframe=teal!35!black,title={Before you move on}]
What does \ru{выносливость} mean? (\textquotedblleft{}Stamina, endurance\textquotedblright{}.) Say it once more.
\end{tcolorbox}
